% s-143-p143-m005; original PDF 143, 144
\Aphorism{s-143-p143-m005}{222}
Von einer unglücklichen Liebe muss man sich gelegentlich in der Einsamkeit erholen, von einer glücklichen muss man es noch viel mehr. „Entziehen wir uns, um aufzuathmen, für einen Augenblick der engen Welt, zu der jede Frage nach dem Werth von \Sperr{Personen} den Geist verurtheilt.“ Liebe ist an sich schon der zudringlichste Affect, den es giebt, eine Gesellschaft, die nie loslässt: ist sie gar noch mit Willigkeit des Objects verbunden, also eine glückliche Liebe, so ist es um Einsamkeit und Selbstbesinnung geschehen — die freie Zeit theilt sich zwischen Liebe und Geliebtem. Diese Beiden können mit einander verschworen sein oder einander entgegenwirken, einander verdrängen; Liebe verstimmt vielleicht gegen den geliebten Gegenstand, und dieser kühlt jene ab durch sein Menschlich-Allzumenschliches. Aber gleichgültig, ob sich die Sache mit dem Bilde, die \Foreign{causa externa} mit der \Foreign{idea concomitans} verträgt oder nicht: beide zusammen gönnen dem Liebenden keine freie Stunde, er muss sie sich selber stehlen, „um aufzuathmen“.
\par

